\documentclass[10pt]{article}

\usepackage[utf8]{inputenc}

\usepackage[T1]{fontenc}

\usepackage{amsmath}

\usepackage{amsfonts}

\usepackage{amssymb}

\usepackage[version=4]{mhchem}

\usepackage{stmaryrd}



\begin{document}

уравнение Пуассона-Больцмана и учитывая условие электронейтральности, получают уравнение



\[

\left(-\Delta+l_{D}^{-2}\right) \varphi(x)=0

\]



к-рое и определяет де баевский экранированный по т е нц и а л:



\[

\varphi(x)=\varphi^{9 \mathrm{~K}}(x)=q \frac{\exp \left[-|x| / l_{D}\right]}{4 \pi|x|},

\]



а также искомые значения давления и плотности в системе.



Строгое доказательство Д.-Х. т. для зарядово симметричных систем в пределе, когда $\varepsilon=1 / k T l_{D}$ стремится к нулю, дано в работе [2]. Сушествование дебаевского экранирования для малых $\varepsilon \neq 0$ было установлено в работах [3]-[4]. Обобщение на случай квантовых кулоновских систем сделано в работе [5].



Jum.: [1] D e by e P., H u c kel E., "Phys. Z.», 1923, Bd 24, № 9, S. 185; [12 Ke n nedy T., "Comm. Math. Phys.», 1983, v. 92, № 2, p. 269; [3] Brydges D., там жe, 1978, v. 58, № 3, p. 313; [4] B rydge s D., Fe de rbush P., там жe, 1980, v. 73, № 3, p. 197; [5] Fon t a i n e J.P., там же, 1986, v. 103, № 2, p. $24 . \quad$ А.Л. Ребенко. ДЕ-БРОЙЛЯ ПЛИНА т е пло во й вол ны - см. АктивHOCmb.



ДЕЙСТВИЕ - фундаментальная физическая величина, задание которой как функции переменных, описывающих состояние системы, полностью определяет динамику системы. Исторически понятие Д. было введено в механике голономных систем (систем со связями, не зависящими от скоростей). Д. $S$ для промежутка времени $\left(t_{1}, t_{2}\right)$ определяется как



\[

S=\int_{t_{1}}^{t_{2}} L\left(q_{i}, \dot{q}_{i}, t\right) d t

\]



где $L=T-U-$ Лагранжа функция, зависящая от описывающих состояние системы обобщенных координат $q_{i}$ и скоростей $\dot{q}_{i}=d q_{i} / d t(i=1, \ldots, n ; n-$ число степеней свободы) и, возможно, времени $t$. При этом кинетич. энергия $T$ квадратична по скоростям, а потенциальная $U$ не зависит от них. Исходными считались уравнения Ньютона, а оправданием для введения понятия Д. служило наблюдение, что эти уравнения получаются как уравнения Эйлера-Лагранжа в вариационном наименьшего действия принципе: $\delta S=0$ при независимых вариациях $\delta q(t)$ с условием $\delta q\left(t_{1}\right)=\delta q\left(t_{2}\right)=0$ на границе.



Уравнениям Лагранжа



\[

\frac{\partial L}{\partial q_{i}}-\frac{d}{d t} \frac{\partial L}{\partial \dot{q}_{i}}=0

\]



эквивалентны Гамильтона уравнения, получающиеся из требования $\delta S=0$ для Д. в эквивалентной (1) форме



\[

S=\int_{t_{1}}^{t_{2}}\left[\sum_{i} p_{i} \dot{q}_{i}-H\left(p_{i}, q_{i}, t\right)\right] d t

\]



при независимых вариациях $\delta q_{i}(t)$ и $\delta p_{i}(t)$ (здесь $H-\Gamma a-$ мильтона функция, $p_{i}$ - обобщенные импульсы). Система обыкновенных дифференциальных уравнений Гамильтона



\[

\dot{q}_{i}=\partial H / \partial p_{i}, p_{i}=-\partial H / \partial q_{i}

\]



служит характеристич. системой для уравнения Гамильтона - Якоби



\[

\frac{\partial S}{\partial t}=-H\left(\frac{\partial S}{\partial q_{i}}, q_{i}, t\right)

\]



к-рое является нелинейным уравнением с частными производными, а интегральные кривые уравнений Гамильтона характеристиками уравнения (3). Д. есть полный интеграл уравнения (3),



\[

S=S\left(\alpha_{i}, q_{i}, t\right)+\alpha_{n+1},

\]



зависящий от $n+1$ произвольных постоянных $\alpha_{k}$, и является производящей функцией канонич. преобразования от переменных $p_{i}, q_{i}$ к новым переменным $P_{i}=\alpha_{i}, Q_{i}=\partial S / \partial \alpha_{i}$. Новая функция Гамильтона $H\left(P_{i}, Q_{i}, t\right)$ тождественно обращается в 0 , вследствие чего новые переменные $P, Q$ постоянны (и выражаются через начальные данные). Тем самым знание полного интеграла (3) сводит задачу интегрирования уравнений движения к разрешению относительно $q_{i}$ алгебраич. уравнений



\[

Q_{i}=\partial S\left(P_{j}, q_{j}, t\right) / \partial P_{i}

\]



В современной теоретич. физике Д. рассматривается как основная фундаментальная величина при формулировке любой теории, особенно полевой, а динамич. уравнения выводятся из вариационных принципов классической механики. Задача построения теории формулируется как задача выбора обобщенных координат и скоростей, описывающих состояние системы, и вида функции Лагранжа, зависящей от них. Значение понятия Д. возрастает для полевых систем еще и потому, что важнейшие для них принципы инвариантности формулируются наиболее удобно и компактно как инвариантность Д. (см. Лагранжев формализм, Лагранжиан); в ряде случаев соображения инвариантности почти полностью определяют теорию. Напр., электродинамикой без источников называется теория, где в качестве координат выбирают 4-потенциал $A_{\mu}(x)$, а требования релятивистской и калибровочной инвариантности и линейности уравнений поля фиксируют Д. в виде



\[

S=\int\left(\frac{\partial A_{\mu}}{\partial x_{v}}-\frac{\partial A_{v}}{\partial x_{\mu}}\right)^{2} d x

\]



где $x=(\boldsymbol{x}, t) \equiv\left\{x_{\mu}\right\}$ - точка пространства-времени. Кроме того, благодаря Нетер теореме инвариантность Д. относительно каждой однопараметрич. группы преобразований влечет за собой закон сохранения одной, явно строящейся по функции Лагранжа (или функции Гамильтона) физич. величины.



Не менее фундаментальна роль Д. в квантовой теории, где состояния системы описываются векторами гильбертова пространства, а динамич. перменным отвечают операторы. Если базис пространства одномерной системы образован собственными векторами $|q\rangle$ оператора координаты, то стандартному постулату квантования эквивалентно определение амплитуды перехода $\left\langle q_{2}\left(t_{2}\right) \mid q_{1}\left(t_{1}\right)\right\rangle$ из состояния с координатой $q_{1}$ в момент $t_{1}$ в состояние с координатой $q_{2}$ в момент $t_{2}$ как функционального интеграла



\[

\left\langle q_{2}\left(t_{2}\right) \mid q_{1}\left(t_{1}\right)\right\rangle=\int \prod d q(t) \exp \left(-i / \hbar \int_{t_{1}}^{t_{2}} L(q, t)\right) d t,

\]



где П (знак умножения) показывает, что интегрирование экспоненты от класич. Д. ведется по всем возможным траекториям, начинающимся в $q_{1}$ в момент $t_{1}$ и кончающимся в $q_{2}$ в момент $t_{2}$. Такая функциональная формулировка особенно удобна для квантовой теории поля: она позволяет ясно следить за инвариантностью на всех этапах, в частности в процедуре перенормировки. Наконец, функциональная формулировка (4) проясняет переход к классич. теории: в квазиклассич. пределе $\hbar \rightarrow 0$, где фазы $S / \hbar$ велики, основной вклад в интеграл дает область, где $S$ стационарно, то есть $\delta S=0$ при вариации траекторий. Таким образом, принцип наименьшего действия для классич. траекторий оказывается следствием квантовой динамики в квазиклассич. пределе. В определенном смысле Д. «более важно» для квантовой теории, чем для классической: квантовую динамику определяют все возможные траектории, а классическую - лишь экстремали.



Лum.: [1] Л а н д а у Л. Д., Л и фшиц Е. М., Теория поля, 6 изд., М., 1973; [2] и х же, Механика, 3 изд., М., 1973; [3] Д и р ак П., Принципы квантовой механики, пер. с англ., 2 изд., М., 1979; [4] Медведе в Б. В., Начала теоретической физики, М., 1977; [5] Р амон П., Теория поля, пер. с англ., М., 1984. В. П. Павлов. ДЕЙСТВИЕ алге $б$ р Ы Л и - см. Ли грynna.



ДЕЙСТВИЕ В д о л в п у т и $\gamma:[0,1] \rightarrow M$ на симплектическом многообразии $М$ с симплектической структурой $\omega=d \lambda-$ функционал, значение которого на $\gamma$ равно



\[

S(\gamma)=\int_{\gamma}(\lambda-H) d t

\]





\end{document}